% ===== 原文 pp.35–36 =====
% [原文 p.35]
% 说明：图 12 在原文 p.35 的顶端，紧随 p.34 末句「…Therefore the gauge group」之后；
% 这里按原文位置放在本片开头。
\begin{center}
\begin{minipage}{\textwidth}
\centering
\includegraphics[width=0.5\textwidth]{fig12.png}\par\vspace{2pt}
\figtitle{12}{两种不同分解下的六穿刺球面}
\end{minipage}
\end{center}

不过，左图与右图的分解相当不同。例如，右图中来自中心那个球面的一个三基表示，其三个
$\SU(2)$ 对称性全都耦合到动力学规范场。但我们也知道，这两种分解可以连续地形变到彼此。
从四维拉格朗日量的角度看，这代表一种 S-对偶。

\subsubsection{对应关系的其他推导}

在第 5.1 节与第 5.2 节中，我们计算了 $\cS(C)$ 在 $S^1\times S^3$ 上以及在 $S^4$ 上的
配分函数，发现它分别由 $C$ 上的二维 $q$-形变 $\SU(2)$ 杨--米尔斯理论和 $C$ 上的
Liouville 理论给出。考虑到当配分函数不依赖于尺寸时成立的关系
$\cS(X_4)(C_2)=\cS(C_2)(X_4)$，我们现在可以断定
\begin{equation}
\cS(S^1\times S^3)=\text{二维 }q\text{-形变 }\SU(2)\text{ 杨--米尔斯理论},
\tag{5.25}
\end{equation}
\begin{equation}
\cS(S^4)=\text{Liouville 理论}.
\tag{5.26}
\end{equation}

% [Tutor Note: p.35, 关系 $\cS(X_4)(C_2)=\cS(C_2)(X_4)$ 是 (5.25)、(5.26) 的全部内容：
%  先在 $C$ 上紧致化、再在 $X_4$ 上算配分函数，与先在 $X_4$ 上紧致化、再在 $C$ 上算，
%  给出同一个理论。这个「先约化哪一边都行」之所以合法，靠的正是 §5.3.1 指出的性质：
%  部分拓扑扭转之后，可计算的四维物理只依赖 $C$ 的总面积与复结构，而不依赖度规细节，
%  因此 $C$ 的尺寸可以任意收缩——原文「when the partition function does not depend on
%  the size」就是这一前提。]

图 13 给出图示。

\begin{center}
\begin{minipage}{\textwidth}
\centering
\includegraphics[width=0.55\textwidth]{fig13.png}\par\vspace{2pt}
\figtitle{13}{六维理论在 $S^4\times$ Riemann 曲面上，及其维数约化}
\end{minipage}
\end{center}

这两个等式很不寻常：在物理学文献中，等式通常是在数之间成立的。而这里的等式却是在量子场论
之间成立！作者相信，这表明我们在更好地理解一般量子场论的道路上正稳步前进。

% [原文 p.36]

我们绕了相当大的弯子，才得到关于 $\cS(X_4)$ 的两个等式 (5.25)、(5.26)：先研究四维理论
$\cS(C)$，再用局域化计算它们的配分函数，最后把结果与已知的二维理论相对照。有没有更直接的
途径得到这些关系？

障碍在于我们不知道六维理论的拉格朗日量。一个可行的实用办法是先做 $S^1$ 约化。当
$X_4=S^1\times S^3$ 时，我们显然有一个 $S^1$。先把六维理论在这个 $S^1$ 上紧致化，剩下的
就是 $S^3$ 上的五维超杨--米尔斯理论。它的路径积分可以局域化到沿 $S^3$ 的常数模，本质上
给出余下两个方向上的 $\SU(2)$ 杨--米尔斯理论。$S^3$ 上的单圈涨落随后给出把它变成
$q$-形变杨--米尔斯理论所需的修正。这一分析在 [36] 中完成。

% [Tutor Note: p.36, 「局域化到沿 $S^3$ 的常数模」这一步可以这样理解：五维超杨--米尔斯在
%  $S^3$ 上的路径积分被局域化压到「$S^3$ 上处处相同的常数规范场」这一鞍点轨迹上；这些
%  常数模正是二维理论里「沿 $S^1$ 的和乐」自由度（也就是 §5.1 里 $Z$ 所依赖的三个
%  $(a_i,1/a_i)$），于是 $S^3$ 方向不再留下动力学自由度，剩下的就是二维杨--米尔斯；
%  $S^3$ 上的单圈涨落则给出把普通杨--米尔斯修正成 $q$-形变杨--米尔斯的那一步。]

当 $X_4=S^4$ 时，我们并不直接看到一个 $S^1$。但 $S^4$ 可以看作 $S^3$ 在一个线段上的
纤维化，并且在两端 $S^3$ 进一步退化。而 $S^3$ 又是 $S^1$ 在 $S^2$ 上的纤维化。于是我们
可以先在这个 $S^1$ 上约化系统。此时系统是：线段上 $S^2$ 纤维化上的五维 $\SU(2)$ 超
杨--米尔斯理论，两端带有某种古怪的边界条件。$S^2$ 上的五维 $\SU(2)$ 超杨--米尔斯理论可以
局域化，给出余下三维中的 $\mathrm{SL}(2)$ Chern--Simons 理论，而线段端点的边界条件使得它
紧致化到二维时给出 Liouville 理论。细节可在文献 [37] 中找到。

% [Tutor Note: p.36, 「$S^4$ 是 $S^3$ 在线段上的纤维化」的具体形式：把
%  $S^4\subset\R^5$ 按 $x_0\in[-1,1]$ 分层，$x_0$ 处的纤维就是半径 $\sqrt{1-x_0^2}$ 的
%  $S^3$，它在两端 $x_0=\pm1$ 处退化为点；而每个 $S^3$ 又是 $S^1$ 在 $S^2$ 上的 Hopf
%  纤维化。所以「先在这个 $S^1$ 上约化」是把纤维方向的圆约掉，两步之后剩下的基底正是
%  线段 $\times$ $S^2$ 的二维部分。]

% [Tutor Note: p.36, 这里出现的是 $\mathrm{SL}(2)$ 而不是紧致的 $\SU(2)$：五维 $\SU(2)$
%  超杨--米尔斯在 $S^2$ 上局域化后给出的三维 Chern--Simons 理论带的是非紧致/复化的规范群，
%  它与 Liouville 理论的联系在规范/引力对偶的文献中很常见；原文把细节推给了 [37]，
%  本译本不再展开。]
